\documentclass[12pt,a4paper]{article}
\usepackage[margin=2.5cm]{geometry}
\usepackage{booktabs}
\usepackage{longtable}
\usepackage{array}
\usepackage{xcolor}
\usepackage{hyperref}
\usepackage{listings}
\usepackage{amsmath}
\usepackage{graphicx}
\usepackage{parskip}

\hypersetup{colorlinks=true, linkcolor=blue, urlcolor=blue}

\lstdefinestyle{code}{
  basicstyle=\ttfamily\small,
  breaklines=true,
  frame=single,
  backgroundcolor=\color{gray!10},
  columns=flexible,
  numbers=left,
  numberstyle=\tiny\color{gray},
  numbersep=5pt,
}

\lstdefinestyle{decaf}{
  basicstyle=\ttfamily\small,
  breaklines=true,
  frame=single,
  backgroundcolor=\color{blue!5},
  columns=flexible,
  keywordstyle=\color{blue}\bfseries,
  morekeywords={class,interface,extends,implements,void,int,double,bool,string,
                if,else,while,for,return,break,null,this,new,New,NewArray,
                Print,ReadInteger,ReadLine,true,false},
}

\title{\textbf{Decaf Mini-Compiler}\\
       \large CS-471L Compiler Construction Lab\\
       University of Engineering \& Technology, Lahore\\
       Spring 2026}
\author{Compiler Project Team}
\date{}

\begin{document}
\maketitle
\tableofcontents
\newpage

% ------------------------------------------------------------------
\section{Introduction}
% ------------------------------------------------------------------

This project implements a front-end mini-compiler for the \textbf{Decaf}
programming language -- a statically typed, object-oriented language designed
specifically for compiler construction courses. Decaf borrows its syntax from
Java but strips away generics, packages, and complex I/O, leaving a tractable
yet expressive target language.

The compiler is implemented in Python 3.10+ and comprises six cooperating
modules as summarised in Table~\ref{tab:modules}.

\begin{table}[h]
\centering
\begin{tabular}{lll}
\toprule
\textbf{Module} & \textbf{File} & \textbf{Role} \\
\midrule
Lexical Analyzer  & \texttt{src/lexer.py}        & Double-buffered tokeniser \\
Recursive Descent & \texttt{src/rd\_parser.py}   & Hand-coded LL parser; builds AST \\
LL(1) Parser      & \texttt{src/ll\_parser.py}   & Table-driven predictive parser \\
SLR(1) Parser     & \texttt{src/lr\_parser.py}   & Bottom-up shift-reduce parser \\
Symbol Table      & \texttt{src/symbol\_table.py}& Scoped hash-based symbol store \\
Error Handler     & \texttt{src/error\_handler.py}& Diagnostics and recovery \\
\bottomrule
\end{tabular}
\caption{Compiler modules and their roles.}
\label{tab:modules}
\end{table}

% ------------------------------------------------------------------
\section{Language Overview: Decaf}
% ------------------------------------------------------------------

A Decaf program is a sequence of top-level declarations.

\subsection{Type System}

\begin{table}[h]
\centering
\begin{tabular}{ll}
\toprule
\textbf{Keyword} & \textbf{Meaning} \\
\midrule
\texttt{int}     & 32-bit signed integer \\
\texttt{double}  & IEEE 754 double precision float \\
\texttt{bool}    & Boolean (\texttt{true} / \texttt{false}) \\
\texttt{string}  & Immutable character string \\
\texttt{void}    & Return type for procedures \\
\textit{ClassName} & Instance of a user-defined class \\
\texttt{T[]}     & One-dimensional array of type T \\
\bottomrule
\end{tabular}
\caption{Decaf type system.}
\end{table}

\subsection{Classes}

\begin{lstlisting}[style=decaf]
class Dog extends Animal implements Runnable {
    string name;
    void bark() { Print("Woof!"); }
}
\end{lstlisting}

Classes support single inheritance (\texttt{extends}) and multiple interface
implementation (\texttt{implements}).

\subsection{Built-in Operations}

\begin{table}[h]
\centering
\begin{tabular}{lll}
\toprule
\textbf{Operation} & \textbf{Return} & \textbf{Description} \\
\midrule
\texttt{Print(e, ...)}      & void    & Print expressions to stdout \\
\texttt{ReadInteger()}      & int     & Read integer from stdin \\
\texttt{ReadLine()}         & string  & Read line from stdin \\
\texttt{New(T)}             & T       & Allocate new class instance \\
\texttt{NewArray(n, T)}     & T[]     & Allocate array of n elements \\
\bottomrule
\end{tabular}
\caption{Decaf built-in operations.}
\end{table}

% ------------------------------------------------------------------
\section{Module 1: Lexical Analyzer}
% ------------------------------------------------------------------

\subsection{Double-Buffering}

The \texttt{DoubleBuffer} class maintains two buffers of 512 characters each.
When the forward pointer reaches the end of one buffer half, the other half is
loaded from the source. This allows efficient character-at-a-time scanning with
O(1) amortised cost per character, and guarantees that tokens spanning a buffer
boundary are handled correctly.

Key operations:
\begin{itemize}
  \item \texttt{next\_char()}: advance the forward pointer; return character.
  \item \texttt{retract()}: move forward pointer back by one.
  \item \texttt{peek()}: look ahead without advancing.
  \item \texttt{start\_lexeme()} / \texttt{get\_lexeme()}: mark and retrieve the
        current token text.
\end{itemize}

\subsection{Token Categories}

\begin{table}[h]
\centering
\begin{tabular}{lll}
\toprule
\textbf{Category} & \textbf{Examples} & \textbf{Notes} \\
\midrule
Keywords  & \texttt{void}, \texttt{class}, \texttt{if} & 24 keywords total \\
Identifiers & \texttt{myVar}, \texttt{Animal} & Max 31 characters \\
Int literals & \texttt{42}, \texttt{0xFF} & Decimal and hexadecimal \\
Double literals & \texttt{3.14}, \texttt{1.0e-5} & With optional exponent \\
String literals & \texttt{"hello"} & Double-quoted \\
Bool literals & \texttt{true}, \texttt{false} & Stored as \texttt{BOOL\_CONST} \\
Operators & \texttt{+}, \texttt{!=}, \texttt{\&\&} & All two-char ops recognised \\
\bottomrule
\end{tabular}
\caption{Lexer token categories.}
\end{table}

% ------------------------------------------------------------------
\section{Module 2: Recursive Descent Parser}
% ------------------------------------------------------------------

The recursive descent parser (RDP) has one Python method per non-terminal.
It builds a dict-based AST and populates the symbol table during parsing.

\subsection{AST Node Structure}

\begin{lstlisting}[style=code,language=Python]
{
    'type':     'FuncDecl',       # node type
    'value':    'main',           # optional: name, operator, literal
    'line':     2,                # source line
    'children': [ret_node, ...]   # child nodes
}
\end{lstlisting}

\subsection{Error Recovery}

Panic-mode recovery is used at the statement level: on a syntax error, tokens
are discarded until a token in the FOLLOW set of the current non-terminal is
seen. This allows reporting multiple independent errors in one pass.

% ------------------------------------------------------------------
\section{Module 3: LL(1) Non-Recursive Predictive Parser}
% ------------------------------------------------------------------

The LL(1) parser uses an explicit stack and a precomputed parse table $M[A, a]$.

\subsection{Parse Table Construction}

For each production $A \to \alpha$:
\begin{enumerate}
  \item For each $t \in \text{FIRST}(\alpha) \setminus \{\varepsilon\}$:
        $M[A, t] = (A \to \alpha)$.
  \item If $\varepsilon \in \text{FIRST}(\alpha)$: for each $t \in \text{FOLLOW}(A)$:
        $M[A, t] = (A \to \alpha)$.
\end{enumerate}

\subsection{Parsing Algorithm}

\begin{lstlisting}[style=code,language=Python]
stack = ['$', 'Program']
token = lexer.next_token()
while stack:
    top = stack[-1]
    sym = token.grammar_symbol()
    if top == '$' and sym == '$':
        break  # accept
    elif top == sym:  # terminal match
        stack.pop(); token = lexer.next_token()
    elif top in table and sym in table[top]:
        production = table[top][sym]
        stack.pop()
        for s in reversed(production.rhs):
            if s != 'epsilon':
                stack.append(s)
    else:
        error(); recover(FOLLOW[top])
\end{lstlisting}

% ------------------------------------------------------------------
\section{Module 4: SLR(1) LR Parser}
% ------------------------------------------------------------------

The SLR(1) parser is a bottom-up, shift-reduce parser producing 229 states
for the Decaf grammar.

\subsection{LR(0) Item Construction}

Closure and goto operations build the canonical LR(0) item collection. An LR(0)
item $[A \to \alpha \bullet \beta]$ shows how far parsing has progressed through
a production.

\subsection{Conflict Resolution}

\begin{table}[h]
\centering
\begin{tabular}{lllll}
\toprule
\textbf{\#} & \textbf{State} & \textbf{Token} & \textbf{Type} & \textbf{Resolution} \\
\midrule
1 & 218 & \texttt{else} & S/R & Prefer shift (nearest if) \\
2 & 108 & \texttt{[}   & S/R & Prefer reduce \\
3 & 108 & \texttt{)}   & R/R & Use first production \\
4 & 1   & \texttt{id}  & S/R & Prefer reduce \\
\bottomrule
\end{tabular}
\caption{SLR(1) conflict resolutions.}
\end{table}

% ------------------------------------------------------------------
\section{Module 5: Symbol Table Manager}
% ------------------------------------------------------------------

The symbol table uses a stack of Python dicts to implement nested lexical scoping.

\subsection{Polynomial Rolling Hash}

\[
  h(s) = \sum_{i=0}^{|s|-1} \text{ord}(s[i]) \cdot 31^i \bmod 211
\]

The prime modulus 211 minimises hash clustering. Each symbol entry stores its
hash value for educational display in the symbol table dump.

\subsection{Scope Stack}

\begin{itemize}
  \item \texttt{scopes[0]}: global (functions, classes, interfaces)
  \item \texttt{scopes[1]}: function parameters
  \item \texttt{scopes[2]}: function body
  \item \texttt{scopes[3+]}: nested block bodies
\end{itemize}

\texttt{lookup(name)} searches from innermost to outermost, implementing the
standard lexical scoping rule.

% ------------------------------------------------------------------
\section{Module 6: Error Handler}
% ------------------------------------------------------------------

\subsection{Error Classification}

\begin{table}[h]
\centering
\begin{tabular}{lll}
\toprule
\textbf{Kind} & \textbf{Raised By} & \textbf{Examples} \\
\midrule
LEXICAL  & Lexer      & Unexpected \texttt{@}, unterminated string \\
SYNTAX   & All parsers& Expected \texttt{;}, found \texttt{\}} \\
SEMANTIC & RD Parser  & Undefined identifier, duplicate declaration \\
\bottomrule
\end{tabular}
\caption{Error classifications.}
\end{table}

\subsection{Panic-Mode Recovery}

After a syntax error, input tokens are discarded until a token in the FOLLOW set
of the offending non-terminal is seen. This allows parsing to continue and report
multiple independent errors.

% ------------------------------------------------------------------
\section{Grammar Transformations}
% ------------------------------------------------------------------

\subsection{Left-Recursion Removal}

Original Decaf uses left-recursive type rules and expression rules. For example:
\begin{align*}
  \textit{Type} &\to \textit{Type}\;\texttt{[]}  \quad\text{(original, left-recursive)}\\
  \textit{Type} &\to \texttt{int}\;\textit{TypeSuffix}
    \mid \ldots\quad\text{(after transformation)}\\
  \textit{TypeSuffix} &\to \texttt{[]}\;\textit{TypeSuffix} \mid \varepsilon
\end{align*}

\subsection{Grammar Statistics}

\begin{center}
\begin{tabular}{lr}
\toprule
Property & Value \\
\midrule
Non-terminals           & 55 \\
Productions             & 121 \\
Terminals               & 40 \\
LL(1) table cells (non-empty) & $\approx$400 \\
SLR(1) states           & 229 \\
SLR(1) conflicts (resolved) & 4 \\
\bottomrule
\end{tabular}
\end{center}

% ------------------------------------------------------------------
\section{Sample Programs}
% ------------------------------------------------------------------

\subsection{test1\_valid.decaf}

\begin{lstlisting}[style=decaf]
void main() {
    int x;
    x = 5;
    int y;
    y = 10;
    int z;
    z = x + y;
    Print(z);
}

int add(int a, int b) {
    return a + b;
}
\end{lstlisting}

All three parsers accept this program. The RD parser builds a two-function AST;
the symbol table records \texttt{main} and \texttt{add} at global scope with their
local variables at scope level 1.

\subsection{test2\_class.decaf}

Tests class declaration, inheritance (\texttt{extends}), member functions, and
\texttt{New()} object allocation. The symbol table correctly creates a class scope
for each class body.

% ------------------------------------------------------------------
\section{Limitations and Future Work}
% ------------------------------------------------------------------

\textbf{Limitations:}
\begin{itemize}
  \item No semantic type checking (types recorded but not verified).
  \item No code generation (front-end only).
  \item Single-file compilation.
\end{itemize}

\textbf{Future work:}
\begin{itemize}
  \item Semantic analysis phase: type checking, overriding validation.
  \item Intermediate code generation (three-address code or LLVM IR).
  \item Optimisation passes: constant folding, dead code elimination.
  \item Language server protocol (LSP) integration.
\end{itemize}

% ------------------------------------------------------------------
\section{References}
% ------------------------------------------------------------------

\begin{enumerate}
  \item Aho, A.~V., Lam, M.~S., Sethi, R., \& Ullman, J.~D. (2006).
        \textit{Compilers: Principles, Techniques, and Tools} (2nd~ed.).
        Addison-Wesley.
  \item Lam, M.~S. et~al. (2003). \textit{CS143 Decaf Language Specification}.
        Stanford University.
  \item Cooper, K.~D., \& Torczon, L. (2011).
        \textit{Engineering a Compiler} (2nd~ed.). Morgan Kaufmann.
  \item Python Software Foundation. \textit{Python 3 Documentation}.
        \url{https://docs.python.org/3/}
\end{enumerate}

\end{document}
